\section{Related Work}
\label{ch:related}

The decision axis of this paper has several close relatives in the literature, and the incremental contribution of this paper can only be stated clearly by sorting out the relations branch by branch.

The first branch is amortized inference. Amortized schemes in variational inference replace per-instance optimization with a recognition network whose parameters are fixed at inference time and whose values are determined by the corpus; this is exactly the inference-theoretic form of Mode B, and recent developments in simulation-based inference and amortized Bayesian workflows continue this form \citep{cranmer2020,radev2023}. The amortization gap, namely the gap between the amortized posterior family and the per-instance optimal posterior, has been identified by theoretical and empirical studies as an attributable component of inference suboptimality in VAE-type models \citep{cremer2018}. This gap is isomorphic to the risk gap of Lemma 3 in this paper: the fixed function family of Mode B corresponds to the amortized family, and the instantiated rules of Mode A correspond to the per-instance optimum. The amortization-gap literature therefore provides evidence, independent of the toy instance, for this paper's premise that the risk gap is positive. The approximate-inference operators that keep this gap open in practice, namely loopy belief propagation, structured variational inference, and learned proposals, are classified in Section~\ref{ch:falsify}: they change the cost and the accuracy of inference, not the provenance of the rule.

The second branch is optimization-based per-instance inference and transductive learning. Per-instance optimization solves each input separately, with rule values instantiated by current evidence, and belongs to Mode A. Transductive learning exploits the test points themselves at inference time \citep{vapnik1995} and likewise lies on the Mode A side, but classical transduction still assumes that test points and training points share a distribution, whereas the open-set case is stronger. Transduction is the precedent of Mode A under near-distribution conditions; the novelty of this paper is to establish Mode A as a necessary condition for open inversion tasks rather than an optional strategy.

The third branch is No-Free-Lunch and its family. This branch is the closest relative of this paper and deserves a layer-by-layer comparison. The original form of NFL concerns the uniformly averaged performance of algorithms over \emph{all} problems and concludes that no algorithm dominates universally \citep{wolpert1996}. The load-bearing premise is the uniform average itself, which flattens task structure and therefore says nothing about whether some algorithm is complete on some task class. The companion result of the same author supplies the other half: given an explicit prior or assumption, a priori distinctions between algorithms \emph{do} exist \citep{wolpert1996b}, and the author's later retrospective stresses that the point of NFL lies in the premises rather than in the slogan that no algorithm is better \citep{wolpert2021}. The optimization version pushes the conclusion into search and optimization \citep{wolpert1997}, and the conservation-of-performance reading rewrites ``indistinguishability under uniform averaging'' as ``improving one class of problems must be paid for by degradation on another'' \citep{ho2002}. Statistical learning theory gives the hypothesis-class version of the same fact: learning is impossible on unrestricted hypothesis classes without inductive bias, and PAC learnability requires controlled capacity \citep{shalev2014}; once a task distribution is made explicit, the bias itself can be learned across tasks \citep{baxter2000}. Learnability has still harder boundaries, and the decidability of some learning problems is itself undecidable \citep{bentdavid2019}. Conceptual work systematically distinguishes the variants of NFL and identifies which conclusion depends on which notion of uniformity \citep{sterkenburg2021}, with a review in \citet{adam2019}.

The relation of this paper to this branch can be located in one sentence: this paper is an instantiation of Wolpert's converse on the class of open inversion tasks, and it lowers the negation from the algorithm level to the component level. The comparison follows in Table~\ref{tab:nfl}.

\begin{table}[t]
\caption{This paper compared with the No-Free-Lunch family}
\label{tab:nfl}
\centering
\small
\begin{tabular}{@{}>{\raggedright\arraybackslash}p{0.15\textwidth}>{\raggedright\arraybackslash}p{0.24\textwidth}>{\raggedright\arraybackslash}p{0.24\textwidth}>{\raggedright\arraybackslash}p{0.25\textwidth}@{}}
\hline
\textbf{Dimension} & \textbf{No-Free-Lunch family} & \textbf{Bias necessity in statistical learning} & \textbf{This paper}\\
\hline
Object of negation & Algorithms under uniform averaging over all problems & Learners without inductive bias & Fixed-memory components, namely rule carriers of frozen provenance\\
\hline
Level of negation & Algorithm level & Hypothesis-class level & Component level\\
\hline
Load-bearing premise & A uniform measure over problems & Unlimited capacity of the hypothesis class & Criterion P4, non-freezability of the demand\\
\hline
Average or pointwise & Averaged over all problems & Uniform convergence in-distribution & Uniform pointwise bound inside the task class\\
\hline
Relation to task structure & Flattens task structure & Requires a prior or capacity control & Fixes the task structure P1 to P4\\
\hline
Testability of the premise & The uniform measure is untestable & Capacity is estimable & P4 admits a testable statement\\
\hline
Consequence of the negation & No algorithm dominates universally & No learning without bias & Frozen provenance cannot carry the guarantee; Mode A is the only route left\\
\hline
\end{tabular}
\end{table}

Three remarks. First, this paper does not contradict NFL but is its converse: NFL says that without a prior there is no distinction, whereas this paper says that given the task structure P1 to P4 and the dichotomy of value sources, the distinction not only exists but falls exactly on the side of the modes. Second, the levels differ: NFL and capacity theory operate at the level of algorithms or hypothesis classes, while the negation of this paper is at the component level, namely that inside one task class, no carrier whose rules are borne by frozen values can satisfy a uniform bound; this is independent of whether some algorithm class is globally optimal. Third, the nature of the premise differs: the load-bearing premise of NFL is a uniform measure that cannot be tested in real settings, whereas the load-bearing premise P4 of this paper admits a testable statement with two separately addressable components, the barrier condition by comparing the demand across scales and the realization level by measuring the frequency of fine-scale demand reversals, as Section~\ref{sec:p4} sets out; correspondingly, the force of the theorems here equals exactly the force of P4, in the sharper form that P4 is a sufficient condition implying a bound rather than a restatement of it, see the logical anatomy in Section~\ref{ch:falsify}.

One possible objection deserves a direct answer: given assumptions, a distinction exists, so is the specialization not trivial. The significance of the specialization lies not in the existence but in the manner and the products. First, the form of the conclusion changes: the converse of NFL yields a comparative ordering in the average sense, whereas this paper yields a pointwise impossibility inside a task class, since no carrier of frozen provenance can satisfy a uniform bound; this is a step from degree to kind. Second, the invariant being distinguished is invisible within the framework of NFL: uniform averaging sees behavior but not the source of the values, while the distinction of this paper falls exactly on the source, so two systems with identical behavior land on opposite sides of the negation because their value sources differ. Third, the prior of NFL is arbitrary whereas the criterion here is not: P4 admits a testable statement, so the theorem can be challenged empirically, while an untestable uniform measure yields no cumulative conclusion. Fourth, the shape of the conclusion differs: an ordering answers which algorithm to swap, a boundary answers which class of architecture to swap, and the boundary here directly yields design implications and a cost trade-off. Fifth, the load-bearing premise is pinned in the open, so strengthening the evidence for P4, adjusting the task class, and upgrading A0 to certifiability all become cumulative future work; in addition, the algebraic skeleton of the argument is machine-verified, and the specialization turns an existence argument into an auditable object.

The fourth branch is distributional robustness and invariance in the training objective. A large family of methods responds to the open set by changing what is optimized rather than where the inference-time values come from. Distributionally robust optimization replaces the empirical objective with a worst-case objective over an ambiguity set of distributions around the corpus \citep{bental2013,rahimian2022}; invariant risk minimization seeks a representation whose optimal predictor is simultaneously optimal across a set of training environments, so that the invariant part is expected to transfer \citep{arjovsky2019}; group-robust objectives minimize the worst group risk when group labels are available, and their study stresses how much regularization is needed for such minima to generalize \citep{sagawa2020}; and domain generalization trains on several source domains with the explicit aim of transferring to an unseen one \citep{zhou2022}, with controlled benchmark comparisons reporting that a carefully tuned empirical-risk baseline remains a strong competitor to the specialized methods \citep{gulrajani2021}. A part of this branch also adapts at test time; those methods are classified by the provenance of their update rule, as the sixth branch below sets out.

The decision axis classifies this branch uniformly rather than competing with it. Every intervention named above acts on the training objective or on the representation being fitted, and none changes the source of the numbers that carry the inference-time rule: the ambiguity set of DRO, the environments of IRM, the group structure of group-robust objectives, and the source domains of domain generalization are all given by the corpus, and the predictor that results is a function with parameters fixed at inference time, hence Mode B by the source criterion, however it was selected. This is not a defect of those methods but the consequence of the question each of them answers: they ask how to select among frozen values, using the corpus, so that the selected values behave better on inputs the corpus does not cover, whereas the axis asks whether any frozen values can carry the guarantee. A robust objective can therefore improve the weak claim, which is where the two literatures agree, and it cannot supply the positive claim, which is what Criterion P4 excludes. The sharpest way to see the difference is that a DRO worst case is defined by an ambiguity set fixed before inference, whereas the open set of this paper is bounded by no set fixed before inference; whether an ambiguity set contains the deployment condition is an assumption of the same family as agreement between the training and test distributions, and it fails in the same way when the corpus does not span the operating conditions. Two boundaries of this branch are worth stating. First, invariance is a property of a frozen representation, so invariant-feature learning moves the burden from the predictor to the representation without changing the source of the values. Second, robustness is purchased against a set specified in advance, so the guarantee of a robust method is inherited from its ambiguity set rather than generated from the current evidence, which is exactly the distinction the decision axis draws.

The fifth branch is hypernetworks. Hypernetworks use one network to generate the weights of another \citep{ha2016}. The generator is the literature counterpart of the ``move one level up'' mechanism of the reach section: the generated weights are values produced at inference time, but the generator's own values are frozen parameters given by the corpus, so Mode B is moved up one level and the ultimate source of the rule's values remains the corpus. The generator is itself a rule, and the procedure of Section~\ref{sec:procedure} asks of it the question it asks of any rule: a corpus-fitted generator is Mode B, while an analytic generator, such as weights computed by projective geometry, belongs to Mode A by the implementer clause of Definition~3. Recursion creates no new source of values; it only lengthens the chain of provenance.

The sixth branch is meta-learning and online adaptation. Here the implementations do change values at inference time, so the question of whether they leave the fixed-memory class is exactly the boundary case the axis is built to settle. Meta-learning fits, across a distribution of tasks, either an initialization or an update rule, so that a few steps on a new task produce an adapted predictor \citep{finn2017}, a program whose methodological spectrum is surveyed by \citet{hospedales2021}. Test-time adaptation updates parameters at inference time from the test input itself, typically under a self-supervised objective \citep{sun2020}, and has grown into a systematic methodological spectrum \citep{liang2024}. The two are classified by the provenance of the update rule and not by the fact that an update occurs. In the meta-learned case, the initialization and the update rule are fitted on a task corpus and are frozen when the new task arrives, so the adaptation is carried out by a corpus-borne rule: Mode B moved up one level, exactly as for hypernetworks. In the test-time case the same criterion is applied to whatever produces the new parameters: an update rule learned on a corpus transfers a corpus-borne function to the current input, whereas an update rule given in closed form from the current evidence, such as a Kalman gain or a strength solved from the current observation by Stein's unbiased risk estimation, instantiates its values in the evidence and belongs to Mode A. The distinction is not academic: few-shot meta-learned adaptation and a single per-instance optimization step both update parameters at test time and land on opposite sides of the axis, which is the boundary case settled in Section~\ref{sec:procedure}. Leaving the fixed-memory class is therefore not the same as escaping the negation, since what escapes is the limit instantiated by current evidence rather than the fact that an optimizer ran; an implementation that takes a few steps from a corpus-initialized point, or that learns its initialization itself, has not moved the source of its values. The control-theoretic and online-learning counterparts of this branch, namely gain scheduling, auto-tuning and model-reference adaptive control, and updates along the current loss sequence, recompute the rule at each step and are Mode A by source \citep{shamma1990,astrom1995,hazan2016}; their guarantees, and the gap between them and the present setting, are discussed in Section~\ref{ch:discussion}.

The seventh branch is the parameter-selection tradition of classical regularization. Since Tikhonov's systematic work \citep{tikhonov1977}, regularization theory has known the limitation of fixed strength and has developed parameter-selection methods such as the L-curve \citep{hansen1992} and generalized cross-validation \citep{golub1979}. Lemma 1 of this paper algebraizes this limitation into a contradiction at the level of optimality conditions and further points out that any parameter selection completed offline still produces a frozen parameter table and hence belongs to Mode B. Against offline selection, this tradition also contains Mode A precedents: Stein's unbiased risk estimation estimates the optimal strength from the current noisy observation itself \citep{stein1981}, anisotropic diffusion and the guided filter take position-wise weights as fixed functions of the current evidence computed online \citep{perona1990,he2013}, and empirical Bayes estimates prior hyperparameters from the current data. These implementations share the extrinsic form of Paradigm E while their source lies on the Mode A side; Section~\ref{ch:falsify} will show that the form is not the criterion, the source is. A neighbouring and widely used form of the same combination appears in the hybrid inverse literature: plug-and-play priors let a denoiser carry the prior while an analytic data-consistency step produces the solution \citep{venkatakrishnan2013}, and regularization by denoising writes the prior of a denoiser into an explicit regularizer \citep{romano2017}. Both are composites whose credit decomposes as the reach section requires, a frozen prior contributing a Mode B candidate without a guarantee and an evidence-evaluated data-consistency step contributing the analytic criterion. Data-consistent learned reconstruction combines a frozen prior with an analytic measurement consistency \citep{haltmeier2026}, an existing form of frozen components proposing candidates with analytic constraints carrying the guarantee; systems of this kind show robust performance within and to a certain extent beyond the training distribution on a standard benchmark \citep{leuschner2021}, a success at the level of the weak claim, consistent with the delimitation of this paper and detailed in the composite-path section of Section~\ref{ch:discussion}.

The eighth branch is the mathematical theory of memory. The recent monograph by Buchanan, Pai, Wang, and Ma systematically grounds deep representation learning as learning structured representations of low-dimensional distributions from data and plainly calls such a representation empirical memory \citep{buchanan2026}; its white-box networks interpret frozen weights as implementations of unrolled optimization algorithms. The book points in the same direction as our decision axis: it states in its own closing chapter that open-ended models, however large, belong to a closed world, that open worlds require closed-loop systems with online error correction, and that in its levels of intelligence, regeneration from memory is only the lowest level while autonomous and continuous correction of existing memory from new observations ranks higher. The increment of this paper relative to that programmatic position is to turn the claim that frozen memory belongs to the closed world from a program into a decidable deductive chain with Criteria P1 to P4, with its algebraic skeleton guaranteed by machine verification.

To sum up, the existing literature has recorded the limitations of Mode B at the empirical level, at the algorithmic level, and at the level of the training objective respectively. The contribution of this paper is to raise the limitation to a completeness negation over a task class, to give a unified decision axis, and to guarantee the algebraic skeleton of the argument chain by machine verification.
